% =========================================================
%  Capítulo 6 — Conclusão
% =========================================================

\chapter{Conclusão}
\label{cap:conclusao}

Este trabalho apresentou o desenvolvimento do \textit{Student App}, uma plataforma
\textit{web} gamificada para apoio à gestão de estudos universitários. A solução
integra ferramentas de organização acadêmica -- gerenciamento de disciplinas,
períodos, atividades, avaliações e horários -- com mecânicas de engajamento
baseadas em gamificação, como pontos de experiência, níveis progressivos, moedas
virtuais e sequências de estudo.

O \textit{backend} foi desenvolvido em Java 17 com o \textit{framework} Spring Boot
3.5.6, adotando a Arquitetura Hexagonal como padrão estrutural. Essa escolha
arquitetural demonstrou-se adequada para o contexto do projeto, pois garantiu o
isolamento das regras de negócio da infraestrutura, facilitando os testes unitários
e tornando os componentes de infraestrutura (banco de dados, armazenamento de
arquivos, e-mail) substituíveis de forma independente.

Os objetivos específicos definidos na Seção~\ref{subsec:objetivos-especificos}
foram alcançados:

\begin{itemize}
  \item O módulo de organização acadêmica foi implementado, contemplando todas as
        entidades e seus fluxos de CRUD;

  \item O sistema de gamificação com XP, níveis, moedas e \textit{streak} foi
        integrado às ações do usuário;

  \item As sessões de foco com recompensa de XP proporcional ao tempo dedicado
        foram implementadas;

  \item A Arquitetura Hexagonal foi aplicada de forma consistente em todo o projeto;

  \item A API REST foi documentada via Swagger UI e protegida com JWT;

  \item O sistema de notificações com deduplicação para prazos, limites de faltas
        e conquistas foi implementado.
\end{itemize}

% -----------------------------------------------------------
\section{Limitações}
\label{sec:limitacoes}

[Descreva as limitações do trabalho atual, por exemplo: ausência do frontend,
ausência de testes de integração, falta de avaliação com usuários reais, etc.]

% -----------------------------------------------------------
\section{Trabalhos Futuros}
\label{sec:trabalhos-futuros}

Como desdobramentos naturais deste trabalho, destacam-se as seguintes direções
para trabalhos futuros:

\begin{itemize}
  \item \textbf{Desenvolvimento do \textit{frontend}:} implementação da interface
        do usuário em uma tecnologia moderna como React ou Vue.js, consumindo a
        API REST existente;

  \item \textbf{Avaliação com usuários reais:} realização de um estudo de usabilidade
        e de eficácia da gamificação com estudantes universitários, coletando métricas
        de engajamento e satisfação;

  \item \textbf{Sistema de conquistas (\textit{badges}):} expansão do sistema de
        gamificação com conquistas desbloqueáveis por marcos específicos (e.g., primeira
        semana consecutiva de estudo, 100 sessões de foco concluídas);

  \item \textbf{Ranking e elementos sociais:} introdução de tabelas de classificação
        e funcionalidades colaborativas entre estudantes;

  \item \textbf{Análise de desempenho:} implementação de relatórios e visualizações
        sobre o histórico de estudos, distribuição de tempo por disciplina e evolução
        das notas;

  \item \textbf{Aplicativo móvel:} desenvolvimento de clientes móveis (iOS/Android)
        com suporte a notificações \textit{push};

  \item \textbf{Testes de integração e de carga:} ampliação da suíte de testes com
        testes de integração (\texttt{@DataJpaTest}) e testes de desempenho sob carga.
\end{itemize}
